\section{Interupsi BIOS}
\label{sec:interupsi-bios}

BIOS menyediakan layanan dasar untuk video, \eng{keyboard}, printer,
dan perangkat lain. Layanan ini tersedia sejak komputer menyala,
sebelum DOS dimuat. Bagian ini membahas dua layanan yang paling
banyak dipakai pada program assembly: \code{INT 10h} untuk video dan
\code{INT 16h} untuk \eng{keyboard}.

Pembahasan diringkas dari \cite{gunawan2019}, Bagian 3,
hlm.~19--25, dan Bagian 18, hlm.~190--202, serta \cite{sto2001},
BAB XV, hlm.~78--84.

\subsection{Layanan Video INT 10h}

\code{INT 10h} memakai register \reg{AH} untuk memilih fungsi.
Register lain menyimpan argumen sesuai fungsi yang dipilih.
Fungsi yang sering dipakai dirangkum pada
Tabel~\ref{tab:int10h}.

\begin{table}[H]
    \centering
    \caption{Fungsi \code{INT 10h} yang sering dipakai.}
    \label{tab:int10h}
    \small
    \begin{tabular}{@{}p{1.6cm}p{7.2cm}p{5.0cm}@{}}
        \toprule
        \textbf{AH} & \textbf{Fungsi} & \textbf{Register argumen} \\
        \midrule
        \hex{00} & Menetapkan mode video & \reg{AL} = nomor mode \\
        \addlinespace
        \hex{02} & Menetapkan posisi kursor &
        \reg{DH} = baris, \reg{DL} = kolom, \reg{BH} = halaman \\
        \addlinespace
        \hex{09} & Menulis karakter dan atributnya &
        \reg{AL} = karakter, \reg{BL} = atribut, \reg{CX} = jumlah,
        \reg{BH} = halaman \\
        \addlinespace
        \hex{13} & Menulis string dan atributnya &
        \reg{ES}:\reg{BP} = string, \reg{CX} = panjang,
        \reg{DH}:\reg{DL} = posisi, \reg{BL} = atribut \\
        \bottomrule
    \end{tabular}
\end{table}

Contoh menetapkan mode teks 80 kolom $\times$ 25 baris berwarna:

\begin{lstlisting}[language=MASM, numbers=none]
MOV AH, 00h         ; fungsi set mode video
MOV AL, 03h         ; mode teks 80x25 berwarna
INT 10h
\end{lstlisting}

Contoh memindahkan kursor ke baris 7, kolom 10:

\begin{lstlisting}[language=MASM, numbers=none]
MOV AH, 02h         ; fungsi set posisi kursor
MOV BH, 00h         ; halaman tampilan 0
MOV DH, 07h         ; baris 7
MOV DL, 0Ah         ; kolom 10
INT 10h
\end{lstlisting}

Fungsi \hex{13} lebih praktis untuk menulis satu string sekaligus.
Register \reg{ES}:\reg{BP} menunjuk string, \reg{CX} menyatakan
jumlah karakter, dan \reg{BL} menyimpan atribut warna. Atribut
\hex{0E} berarti kuning pada latar hitam, sedangkan \hex{0C} berarti
merah terang pada latar hitam.

\subsection{Layanan Keyboard INT 16h}

\code{INT 16h} dipakai untuk membaca tombol yang ditekan. Fungsi
utamanya ditunjukkan pada Tabel~\ref{tab:int16h}.

\begin{table}[H]
    \centering
    \caption{Fungsi \code{INT 16h} yang sering dipakai.}
    \label{tab:int16h}
    \small
    \begin{tabular}{@{}p{1.6cm}p{6.8cm}p{5.4cm}@{}}
        \toprule
        \textbf{AH} & \textbf{Fungsi} & \textbf{Hasil} \\
        \midrule
        \hex{00} & Membaca tombol; menunggu sampai ada tombol ditekan &
        \reg{AH} = kode scan, \reg{AL} = kode ASCII \\
        \addlinespace
        \hex{01} & Memeriksa apakah ada tombol ditekan &
        \flag{ZF} = 0 bila ada tombol; kode di \reg{AX} \\
        \addlinespace
        \hex{02} & Membaca status tombol khusus &
        \reg{AL} = status shift \\
        \bottomrule
    \end{tabular}
\end{table}

Fungsi \hex{00} bersifat menunggu (\eng{blocking}): program berhenti
sampai pengguna menekan tombol. Fungsi \hex{01} bersifat tidak
menunggu (\eng{non-blocking}) sehingga program dapat terus berjalan
sambil memeriksa tombol.

\begin{peringatanbox}
\vspace{1.1em}
Layanan \code{INT 16h} membaca tombol pada level BIOS, sedangkan
\code{INT 21h} fungsi \hex{01} atau \hex{0A} membaca masukan melalui
DOS. Keduanya tidak boleh dicampur tanpa alasan yang jelas karena
dapat memengaruhi isi penyangga \eng{keyboard}.
\end{peringatanbox}
